\section{The question that predates the spec}

Organizations have poured budget, headcount, and roadmap into autonomous
software work, and most of them still cannot say whether it paid off. MIT Media
Lab's Project NANDA, surveying roughly three hundred enterprise generative-AI
initiatives, found that about 95\% showed no measurable P\&L return
\citep{nanda2025}. That figure has many causes. One of them is plain: a return
you cannot attribute is a return you cannot claim. The 95\% is, in part, an
attribution failure.

Adoption is not the constraint. Stanford's AI Index put enterprise AI adoption
at 78\% in its 2025 edition and 88\% a year later
\citep{stanford2025, stanford2026}. The tools are in the building. What remains
unanswered is whether the autonomous work they produce is delivering, and for
whom.

Scale is what makes the gap urgent. A single developer reviewing a single
assistant's suggestions is a manageable trust problem. A fleet of agents opening
changes across dozens of repositories, at a rate no human reads line by line, is
not. Trust at that scale has to be granted by policy rather than by inspection,
and policy needs a signal to run on. The governance question here is older than
any protocol. Before a wire format made agent activity legible, leaders already
needed to know which autonomous work had earned more scope and which had not.
The Model Context Protocol revision finalizing on 2026-07-28 does not create that
question. It changes what an organization can answer it with.

The question also already has a named answer in the field, and that answer is
worth taking seriously before improving on it. Earned autonomy, graduated
autonomy, and progressive autonomy all describe the same ladder: an agent starts
under human review, and as it performs, it earns the right to act with less
oversight. Matthias Roder's earned-autonomy gradient formalizes the rungs
\citep{roder}; a growing set of platform vendors ship progressive-autonomy
controls \citep{progressiveautonomy}; the Digital Apprentice line of work models
the same apprenticeship arc \citep{digitalapprentice}; and McKinsey frames agent
trust as decision rights granted and revoked over time \citep{mckinsey}. This
ladder is published across vendors and research, and
it is converging on a shared shape.

Every approach we surveyed conditions autonomy on the same class of signal: how
well the agent performed its task. Extraction accuracy, override rate, tool-call
correctness. That choice of variable is the one thing this paper revisits, because
it is the same variable that lets a clean record hide a bad outcome. This paper
keeps the name earned autonomy and changes only the variable it conditions on:
delivered outcome, not task accuracy. The rest of the argument shows why the
2026-07-28 spec is the first open protocol standard on which that condition can
be enforced per action.
